\documentclass[11pt]{article}
\usepackage{preamble}
\renewcommand{\answer}[1]{}

\begin{document}

\ladderheading{AP Statistics \textperiodcentered\ Unit 2 \textperiodcentered\ Topic 2.10 \textperiodcentered\ B: 2026-11-05 \textperiodcentered\ E: 2026-11-25}{Rare Is Not Never}

\focusbanner{TODAY'S PROBLEM}

Suppose a basketball player has made 70 percent of her free throws over a long career. Tonight she makes 3 of 10 attempts. Let $X$ be the number she would make in 10 attempts if her chance on each shot were still 0.70. The graph shows the binomial model as expected frequencies per 1,000 sets of 10 shots.\par\smallskip \textbf{Claim A:} ``She is in a slump---making only 3 of 10 never happens at her career rate.''\par \textbf{Claim B:} ``With only 10 shots, 3 makes is ordinary chance variation.''\par Under this model, technology gives $P(X\leq3)=0.0106$ for 3 or fewer makes.\par
\begin{center}
\begin{tikzpicture}\begin{axis}[width=0.95\linewidth,height=1.8in,ybar,bar width=7pt,xmin=-0.6,xmax=10.6,ymin=0,ymax=290,xtick={0,1,2,3,4,5,6,7,8,9,10},ytick={0,100,200},xlabel={Made shots},ylabel={Count per 1,000},tick label style={font=\small},label style={font=\small},ymajorgrids]\addplot[fill=calloutblue,draw=dayblue] coordinates {(0,0.006) (1,0.138) (2,1.447) (3,9.002) (4,36.757) (5,102.919) (6,200.121) (7,266.828) (8,233.474) (9,121.061) (10,28.248)};\end{axis}\end{tikzpicture}
\end{center}
\begin{firsttakebox}{FIRST TAKE --- before the video (5 min)}
Which claim, if either, do you trust more? Cite one detail from the result or graph.
\workspace{0.9in}
\end{firsttakebox}

\begin{callout}{\IconBook\ BINS BEFORE BINOMIAL (keep on the page)}{calloutgreen}
A binomial count uses Binary outcomes, Independent trials, a fixed Number of trials, and the Same success chance. Define $X$, $n$, and $p$. For exactly $x$ successes, $P(X=x)=\binom{n}{x}p^x(1-p)^{n-x}$. A cumulative probability adds the included counts. A small probability means unusual under the model, not impossible; it does not identify a cause.
\end{callout}

\newpage
\textbf{Finish after the video:}\par\smallskip

\dokbadge{1}~\textbf{(a)} State what $X$ counts, and give $n$ and $p$.\par
\workspace{0.8in}

\dokbadge{2}~\textbf{(b)} Calculate $P(X=3)$ using the binomial formula. Explain how this event differs from $X\leq3$.\par
\workspace{1.3in}

\dokbadge{3}~\textbf{(c)} Is either claim fully supported? Use the given chance of 3 or fewer makes to explain what each claim gets wrong. Name one reason independence or the same success chance might fail, and one kind of new evidence that would help. Write 3--4 sentences.\par
\workspace{2.0in}

\begin{sentenceframebox}\raggedright\textbf{Frame:} The binomial model is \blank[1in] because \blank[2.5in], although independence requires \blank[1.8in].\end{sentenceframebox}

\begin{sentenceframebox}\raggedright\textbf{Frame:} Since $P(X\leq3)=\blank[0.8in]$, Claim \blank[0.4in] is \blank[0.9in], but \blank[2.2in].\end{sentenceframebox}

\begin{summaryexitbox}{EXIT --- turn this in}
One thing the video changed about my first take: \hrulefill\par
\workspace{0.4in}
\end{summaryexitbox}

\end{document}
